\section{Layered Architecture}
\label{sec:architecture}

This section defines the layered architecture of the diagram compiler, emphasizing the clean separation between kernel, grammars, and composition.

\subsection{The Three Layers}

\ourdiagram[0.95\linewidth]{three-layers}{%
  The three-layer architecture: the Composition layer depends on peer Grammar modules,
  each Grammar depends on the Kernel.
  The Kernel (Scene IR, Backends, Themes, Determinism, Icons, Lint) is grammar-agnostic.}

\ourdiagram[0.95\linewidth]{architecture}{%
  The two-frontend, two-IR compilation pipeline.
  Mermaid DSL (\texttt{.mmd}) and structured IR (YAML/JSON) both feed into a Domain IR,
  which the layout engine transforms into a Scene IR consumed by SVG, PNG, and Skia backends.}

\subsection{Layer 1: The Kernel}
\label{sec:arch-kernel}

The kernel is the shared infrastructure that all grammars depend on. It is grammar-agnostic: it knows nothing about timelines, flows, or graphs.

\subsubsection{Kernel Components}

\begin{description}
  \item[Scene IR (§\ref{sec:scene-ir})] The universal primitive set (Rect, Circle, Line, Path, Text, Group) and effect definitions. The render contract.
  
  \item[Rendering Backends (§\ref{sec:backends})] SVG, PNG (resvg), Skia, PPTX backends. Consume Scene IR, produce output.
  
  \item[Theme System (§\ref{sec:themes})] Theme schema, theme loading, inheritance, property resolution. Grammar-agnostic but grammars may extend.
  
  \item[Determinism Contract (§\ref{sec:determinism})] Rounding rules, sort ordering, no-random, no-clock guarantees.
  
  \item[Deterministic Font Metrics] Embedded font metric tables for layout-critical text measurement.
  
  \item[Icon Registry] Named icons (SVG paths) available to all grammars.
  
  \item[Asset/Image Loader] Resolves and embeds external images as data URIs.
  
  \item[Layout Helpers] Common layout primitives: text measurement, box collision, orthogonal routing.
  
  \item[Lint/Quality-Gate Framework] Validation diagnostics, severity levels, suggestion generation.
  
  \item[Animation Hints (§\ref{sec:animation})] Declarative animation attached to Scene primitives.
\end{description}

\subsubsection{Kernel Contract}

The kernel makes the following guarantees to grammars:

\begin{enumerate}
  \item \textbf{Grammar-agnosticism}: The kernel never imports from, or references, any grammar module.
  
  \item \textbf{Determinism}: All kernel operations are deterministic given the same inputs.
  
  \item \textbf{Backend uniformity}: All backends consume the same Scene IR interface.
  
  \item \textbf{Theme uniformity}: Theme properties are resolved identically for all grammars.
\end{enumerate}

\subsection{Layer 2: Grammars}
\label{sec:arch-grammars}

Grammars are peer modules that sit atop the kernel. Each grammar:

\begin{itemize}
  \item Defines its own Domain IR schema
  \item Implements one or more layout engines (Domain IR → Scene IR)
  \item May extend the theme schema with grammar-specific properties
  \item Exposes validation and compilation entry points
\end{itemize}

\subsubsection{Grammar Independence}

Grammars are \emph{independent} of each other:

\begin{itemize}
  \item The Timeline grammar does not import from the Flow grammar
  \item Adding a new grammar does not modify existing grammars
  \item Each grammar can be versioned and released independently
\end{itemize}

\subsubsection{Grammar-Kernel Dependency}

All grammars depend on the kernel:

\begin{lstlisting}[language=javascript,caption={Grammar dependency structure}]
// packages/timeline/package.json
{
  "dependencies": {
    "@diagram-compiler/kernel": "^1.0.0"
  }
}

// packages/flow/package.json
{
  "dependencies": {
    "@diagram-compiler/kernel": "^1.0.0"
  }
}
\end{lstlisting}

\subsection{Layer 3: Composition}
\label{sec:arch-composition}

The Composition layer sits above grammars. It:

\begin{itemize}
  \item References panels, each specifying a grammar type and Domain IR
  \item Invokes grammar compilation for each panel
  \item Arranges panel Scenes into a unified poster Scene
  \item Adds editorial chrome (headers, footers, dividers)
\end{itemize}

\subsubsection{Composition Dependencies}

Composition depends on grammars (to compile panels) and kernel (for chrome rendering):

\begin{lstlisting}[language=javascript,caption={Composition dependency structure}]
// packages/composition/package.json
{
  "dependencies": {
    "@diagram-compiler/kernel": "^1.0.0",
    "@diagram-compiler/timeline": "^1.0.0",
    "@diagram-compiler/flow": "^1.0.0",
    "@diagram-compiler/comparison": "^1.0.0"
    // ... all available grammars
  }
}
\end{lstlisting}

\subsection{Dependency Rules}

\begin{center}
\small
\begin{tabular}{lcccc}
\toprule
\textbf{Module} & \textbf{→ Kernel} & \textbf{→ Grammars} & \textbf{→ Composition} & \textbf{→ Other Grammars} \\
\midrule
Kernel & — & ✗ & ✗ & ✗ \\
Grammar (any) & ✓ & — & ✗ & ✗ \\
Composition & ✓ & ✓ & — & — \\
\bottomrule
\end{tabular}
\end{center}

\paragraph{Key rule:} The kernel must not import anything grammar-specific. This is enforced by linting.

\subsection{Module Boundaries}

\subsubsection{What Lives in the Kernel}

\begin{itemize}
  \item Scene IR types and serialization
  \item All rendering backends
  \item Theme schema (core properties)
  \item Theme loader and resolver
  \item Icon registry
  \item Font metric tables
  \item Validation framework (DiagnosticSeverity, ValidationResult)
  \item Determinism utilities (rounding, sorting)
  \item Animation hint types
\end{itemize}

\subsubsection{What Lives in Grammars}

\begin{itemize}
  \item Domain IR types (e.g., Activity, Track, Milestone for Timeline)
  \item Domain IR schema (JSON Schema)
  \item Validation rules (semantic validation beyond schema)
  \item Layout engine(s)
  \item Grammar-specific theme extensions
\end{itemize}

\subsubsection{What Lives in Composition}

\begin{itemize}
  \item Composition IR types
  \item Composition IR schema
  \item Panel arrangement algorithms (stack, grid, custom)
  \item Chrome rendering (headers, footers, dividers)
  \item Grammar dispatcher (routes panel to correct grammar)
\end{itemize}

\subsection{API Surface}

Each layer exposes a clean API:

\subsubsection{Kernel API}

\begin{lstlisting}[language=javascript,caption={Kernel public API}]
// Scene IR
export type Scene, DrawingPrimitive, EffectDefinition, AnimationHint;
export function sceneHash(scene: Scene): string;

// Backends
export function renderSvg(scene: Scene): string;
export function renderPng(scene: Scene): Buffer;
export function renderPdf(scene: Scene): Buffer;
export function renderPptx(scene: Scene): Buffer;

// Themes
export function loadTheme(name: string): Theme;
export function resolveTheme(base: Theme, overrides: Partial<Theme>): Theme;

// Validation
export type Diagnostic, ValidationResult;

// Icons
export function getIcon(name: string): SvgPath | undefined;

// Font metrics
export function measureText(text: string, font: FontSpec): TextMetrics;
\end{lstlisting}

\subsubsection{Grammar API}

Each grammar exposes the same interface shape:

\begin{lstlisting}[language=javascript,caption={Grammar public API (Timeline example)}]
// @diagram-compiler/timeline
export type TimelineIR;  // Domain IR type
export const schema: JsonSchema;
export function validate(ir: unknown): ValidationResult;
export function compile(ir: TimelineIR, theme: Theme): Scene;
\end{lstlisting}

\subsubsection{Composition API}

\begin{lstlisting}[language=javascript,caption={Composition public API}]
// @diagram-compiler/composition
export type CompositionIR;
export const schema: JsonSchema;
export function validate(ir: unknown): ValidationResult;
export function compile(ir: CompositionIR, theme: Theme): Scene;
\end{lstlisting}
